\section{Comparison of Steady-State Eigenvalues across Calibrations}
\label{app:additional}

Table~\ref{tab:eigenvalues} summarises the five generalised
eigenvalues of the dynamic system in
equation~\eqref{eq:matrix-form} under the baseline calibration of
Table~\ref{tab:calibration} and the three alternative Taylor-rule
calibrations considered in Section~\ref{sec:robustness}. In every case
three eigenvalues coincide with the AR(1) persistence parameters and
two form a complex-conjugate pair outside the unit circle, confirming
that the Blanchard--Kahn condition holds throughout.

\begin{table}[h]
  \centering
  \caption{Dynamics-matrix eigenvalues under alternative Taylor rules.}
  \label{tab:eigenvalues}
  \begin{tabular}{lcccc}
    \toprule
    & Baseline & Hawkish & Dovish & Strict $\phi_{y}=0$ \\
    & $(\phi_{\pi},\phi_{y})=(1.5,0.125)$ & $(2.0,0.125)$ & $(1.1,0.125)$ & $(1.5,0)$ \\
    \midrule
    $\mu_{1}$ & 0.900 & 0.900 & 0.900 & 0.900 \\
    $\mu_{2}$ & 0.500 & 0.500 & 0.500 & 0.500 \\
    $\mu_{3}$ & 0.500 & 0.500 & 0.500 & 0.500 \\
    $|\mu_{4}|$ & 1.182 & 1.215 & 1.156 & 1.171 \\
    $|\mu_{5}|$ & 1.182 & 1.215 & 1.156 & 1.171 \\
    \midrule
    Determinacy & \checkmark & \checkmark & \checkmark & \checkmark \\
    \bottomrule
  \end{tabular}
\end{table}

The three stable eigenvalues are invariant across calibrations because
they are pinned down by the AR(1) shock persistences. The two unstable
eigenvalues form a complex pair whose modulus grows with the
inflation-reaction coefficient $\phi_{\pi}$, reflecting that a more
aggressive central bank generates faster jumper convergence to the
saddle-path equilibrium.
